\documentclass[12pt]{article}
\usepackage[utf8]{inputenc}
\usepackage[T1]{fontenc}
\usepackage{hyperref}
\title{Kuku-Uwanh Language Resources for K-12 Curriculum Development}
\author{Compiled by Claude Sonnet 5 (claude-sonnet-5) for Andr\'as Kornai}
\date{}
\begin{document}
\maketitle

\section{Introduction}
Kuku-Uwanh (ISO 639-3 uwa; Glottocode kuku1280) is a Paman (Pama-Nyungan) variety of Cape York Peninsula, Queensland. Search results describe it as a dialect of Kugu-Nganhcara (Wik-Nganhcara). Other sources describe Kugu-Nganhcara as having fewer than 40 speakers, with only a few elderly speakers. The reference description supplied for this task calls Kuku-Uwanh extinct. The two statements conflict, and I could not resolve the conflict from search results. Very little dedicated material was found. No Kuku-Uwanh wordlist, primer or textbook turned up. This dossier is short because little exists in what I could find.

\section{Primary Educational Resources}
No Kuku-Uwanh-specific school materials were found. The Queensland Government page on 2020 language preservation projects describes an Aurukun Shire Council language program with Koolkan Aurukun State School and Elders. It covers teaching resources for Wik-Ngathan, Wik-Alken/Wik-Ngatharr and Kugu. Its relevance to Kuku-Uwanh itself is unconfirmed.
\begin{itemize}
\item Preservation projects in Queensland (2020): \url{https://www.qld.gov.au/firstnations/cultural-awareness-heritage-arts/culture-arts/international-decade-indigenous-languages/preservation-projects-queensland/2020}
\end{itemize}

\section{Linguistic and Reference Materials}
\begin{itemize}
\item Wikipedia, Kugu Nganhcara language (background on the wider language): \url{https://en.wikipedia.org/wiki/Kugu_Nganhcara_language}
\item Endangered Languages Project entry: \url{https://endangeredlanguages.com/node/92026}
\item Grambank entry for kuku1280, which cites Smith and Johnson 2000 (that publication was not checked): \url{https://grambank.clld.org/languages/kuku1280}
\item State Library of Queensland Aboriginal and Torres Strait Islander language word lists (dataset overview; a Kuku-Uwanh list was not found in it): \url{https://www.data.qld.gov.au/dataset/slq-aboriginal-languages-word-lists}
\end{itemize}

\section{Religious and Community Resources}
Pama Language Centre works on neighbouring Wik languages (Wik-Alken revitalisation, Wik-Mungkan). No Kuku-Uwanh-specific community resource was found.
\begin{itemize}
\item Pama Language Centre, Wik-Alken activities: \url{https://www.pamacentre.org.au/wik-alken-current-language-activities/}
\end{itemize}

\section{Digital and Audio Resources}
No Kuku-Uwanh audio or digital resources were found in search results.

\section{Curriculum Development Context}
Any curriculum would be a reclamation effort built from archival sources. The next step would be to locate the sources cited by Grambank and by linguistic literature on Kugu-Nganhcara. Work should go through Wik and Kugu community organisations such as the Pama Language Centre and the Aurukun language program. These steps are suggestions and are not confirmed by the search results.
\end{document}